
\subsubsection{2.1 Execution-Feedback Reinforcement Learning for Code Generation}

Reinforcement learning from execution feedback has been studied as an alternative to supervised fine-tuning for improving code generation. CodeRL \citep{Leetal2022} introduces actor-critic training with unit test feedback. PPOCoder \citep{Shojaeeetal2023} applies PPO-based RL to code synthesis and reports improvements on HumanEval and MBPP. TÜLU 3 \citep{Lambertetal2024} uses GRPO with binary pass/fail rewards in a multi-task alignment setting.

All of these works measure performance via pass@1 outcomes. None measures the structural character of policy movement relative to the base model or provides a metric for comparing structural change quality across methods.

\subsubsection{2.2 DPO and Preference-Based Alignment for Code}

Direct Preference Optimization \citep{Rafailovetal2023} reformulates RLHF as a supervised learning problem over preference pairs, avoiding online rollouts. DPO has been applied to code generation [Guo et al., 2024; Luo et al., 2023] typically using execution-oracle labeling (passing solutions preferred over failing). GRPO \citep{Shaoetal2024} reformulates PPO without a critic network using group-relative rewards and became prominent following DeepSeek-R1 [DeepSeek-AI, 2025].

KL divergence appears in both frameworks as a training-time constraint: as the implicit regularization in DPO's objective and as the explicit penalty term β in GRPO. Neither framework uses KL divergence as a post-hoc diagnostic for structural movement analysis. The present work repurposes KL divergence from a training constraint into a normalizer for structural efficiency measurement.

\subsubsection{2.3 AST-Based Code Analysis}

The program analysis literature uses AST-level analysis extensively. The FA-AST taxonomy \citep{Wangetal2020} classifies Python AST nodes to support code clone detection with graph neural networks. ZSS tree edit distance [Zhang and Shasha, 1989] has been applied to code clone detection \citep{Svajlenkoetal2014} and automated program repair \citep{Huaetal2018}. Ding et al. [2024] validate AST edit distance against established code similarity metrics.

These methods are applied to pairs of programs, not to the policy movement of a trained model relative to its base. The composition of FA-AST classification, ZSS edit distance, and KL-matched checkpoint comparison — which is necessary for structural efficiency measurement — has not been proposed or empirically evaluated prior to this work.

\subsubsection{2.4 Policy Analysis and Interpretability in RL Fine-Tuning}

Work on understanding RL fine-tuning behavior [Schulman et al., 2017; Zheng et al., 2023] addresses training stability, reward hacking, and KL budget management. Mechanistic interpretability research \citep{Elhageetal2022} analyzes internal model representations. Neither line of work provides code-level structural analysis of the policy change induced by different alignment methods.
